\section[Cote 83 : la droite des twists et son lieu dégénéré]{Cote 83 : la droite des twists et son lieu dégénéré\protect\verifie[\leanblobmain]{Folder83}}\label{sec:83}

La cote 83 réunit des notes manuscrites sur les algèbres d'Azumaya de rang~$4$ ; sa station~X, datée \q{Avril 86}, couvre les pages~101 à~147. Les pages~92 à~98 avaient construit, pour $(b, c, b', c', t) \in k^5$, une algèbre $M(b, c; b', c'; t)$ libre de base $1, u, v, uv$, avec $u^2 = bu - c$, $v^2 = b'v - c'$ et $\operatorname{Tr}(uv) = t$, et le théorème de la page~98 dit qu'elle est d'Azumaya si et seulement si $Q = t^2 - bb't + b^2c' + b'^2c - 4cc'$ est inversible. Les pages~126 à~136 reprennent la question de façon intrinsèque. Pour deux algèbres quadratiques $H$, $H'$ et $P = H \otimes H'$, le \q{Th. récapitulatif} affirme que les structures d'algèbre de Hamilton--Cayley sur $P$ qui induisent les structures données sur $H \otimes 1$ et $1 \otimes H'$ forment un torseur sous $(L \otimes L')^{\vee}$, $L = H/k$, $L' = H'/k$, identifié à celui des \emph{twists} $T : H \times H' \to k$, avec $T(u, 1) = \operatorname{Tr}u$ et $T(1, v) = \operatorname{Tr}v$ (\lot{83}{7}{130}, \lot{83}{7}{136}). Les pages~137 à~143 étudient le lieu $\widetilde S = \delta_{P,B}^{-1}(0)$ où la structure cesse d'être d'Azumaya. À deux caractères $\varphi$, $\varphi'$ de $H$ et $H'$, elles associent le twist $t_{\varphi,\varphi'}(u, v) = \operatorname{Tr}(v)\varphi(u) + \operatorname{Tr}(u)\varphi'(v) - 2\varphi(u)\varphi'(v)$, vérifient qu'il annule $Q$ et identifient l'algèbre de $\widetilde S$ au produit $H * H'$, engendré par $u * v = \operatorname{Tr}(v)u + \operatorname{Tr}(u)v - 2u \otimes v$ (\lot{83}{7}{138}, \lot{83}{8}{143}). La page~138 écrit le terme constant de $Q$ sous la forme $b^2b' + b'^2b - 4cc'$, et la page~143 donne $\det(u * v) = b^2c + b'^2c' - 4cc'$. La lecture corrige les deux en $b^2c' + b'^2c - 4cc'$.

\noindent\emph{Conventions.} Soit $k$ un anneau commutatif, sans hypothèse sur~$2$. On prend $H = k[u]/(u^2 - bu + c)$ et $H' = k[v]/(v^2 - b'v + c')$ libres, ce qui trivialise $L$ et $L'$, donc $(L \otimes L')^{\vee} = k$. On note $\delta = b^2 - 4c$, $\delta' = b'^2 - 4c'$, et $Q(t) = t^2 - bb't + b^2c' + b'^2c - 4cc'$. Sur $P = k^4$, de base $(1, u, v, uv)$, la multiplication de $M(b, c; b', c'; t)$ est la table de la page~92 : $vu = (t - bb') + b'u + bv - uv$, $v(uv) = -bc' + c'u + tv$, $(uv)u = -b'c + tu + cv$, $(uv)v = -c'u + b'uv$, $(uv)^2 = -cc' + t\,uv$. On pose $\operatorname{Tr}(x) = 2x_0 + bx_1 + b'x_2 + tx_3$ pour $x = x_0 + x_1u + x_2v + x_3uv$, et
\[ \det(x) = x_0^2 + cx_1^2 + c'x_2^2 + cc'x_3^2 + bx_0x_1 + b'x_0x_2 + tx_0x_3 + (bb' - t)x_1x_2 + b'cx_1x_3 + bc'x_2x_3. \]

\begin{enonce}[Pages 81 à 84 et 130 à 136, la droite des structures\piste{83-twist-line-degenerate-locus-is-product}]
\begin{enumerate}
\item Pour tout $t \in k$, $M(b, c; b', c'; t)$ est une $k$-algèbre associative et unitaire. Elle induit les structures données sur $H$ et $H'$ ($u^2 = bu - c$, $v^2 = b'v - c'$, $u \cdot v = uv$), tout élément $x$ y vérifie $x^2 - \operatorname{Tr}(x)x + \det(x) = 0$, $\det$ est multiplicatif et $\operatorname{Tr}(uv) = t$.
\item Réciproquement, soit $A$ une $k$-algèbre dont tout élément vérifie $x^2 - T(x)x + N(x) = 0$, pour une forme linéaire $T$ et une fonction $N$. Soient $u, v \in A$ avec $T(u) = b$, $N(u) = c$, $T(v) = b'$, $N(v) = c'$. Il existe alors $t \in k$ tel que $x \mapsto x_0 + x_1u + x_2v + x_3uv$ soit multiplicative de $M(b, c; b', c'; t)$ dans $A$. Si cette application est injective, $t$ est unique.
\item Deux valeurs distinctes de $t$ donnent deux multiplications distinctes. Les structures forment donc la droite affine des $t$, torseur sous $k = (L \otimes L')^{\vee}$, où $\lambda \in k$ agit par $t \mapsto t + \lambda$.
\end{enumerate}
\end{enonce}

\begin{proof}
(1) Chaque axiome est une identité polynomiale en les coordonnées et en $b, c, b', c', t$, que l'on vérifie coordonnée par coordonnée, et l'associativité aussi : la page~84 l'évite par un argument de prolongement des identités, mais elle se vérifie directement sur la table. (2) Hamilton--Cayley appliqué à $u$ et à $v$ donne $u^2 = bu - c$ et $v^2 = b'v - c'$. Appliqué à $u + v$, avec $T(u + v) = b + b'$, il donne la relation~(6') : $vu = (t - bb') + b'u + bv - uv$, où $t = bb' + c + c' - N(u + v)$. Les autres produits s'en déduisent par associativité. Par exemple $v(uv) = (vu)v = (t - bb')v + b'uv + b(b'v - c') - u(b'v - c') = -bc' + c'u + tv$, et de même pour $u(uv)$, $(uv)u$, $(uv)v$ et $(uv)^2$. On retrouve la table de la page~92, d'où la multiplicativité en développant. Si deux valeurs $t$ et $t'$ conviennent, les images de $v \cdot u$ coïncident ; par injectivité, $vu$ est le même élément de $P$ pour les deux structures, et son coefficient de~$1$, $t - bb' = t' - bb'$, donne $t = t'$. (3) Même calcul sur le coefficient de $1$ dans $v \cdot u$.
\end{proof}

\begin{enonce}[Pages 94, 116 à 120 et 137 à 143, le lieu dégénéré et le produit\piste{83-twist-line-degenerate-locus-is-product}]
\begin{enumerate}
\item Sur $B = k1 + ku + kv$, la forme $\det$ est $x_0^2 + cx_1^2 + c'x_2^2 + bx_0x_1 + b'x_0x_2 + (bb' - t)x_1x_2$. Son demi-discriminant $\delta_{P,B}$ (la moitié du déterminant de la matrice polaire, défini sur tout anneau en rang impair) vaut $-Q(t)$. Le polynôme $Q$ a pour discriminant $\delta\delta'$ et est invariant par $t \mapsto bb' - t$. Le discriminant de la forme trace $\operatorname{Tr}(e_ie_j)$ de $M(b, c; b', c'; t)$ dans la base $(1, u, v, uv)$ vaut $-Q(t)^2$. Enfin $\widetilde S = \operatorname{Spec} k[t]/(Q)$ est libre de rang~$2$ sur~$k$.
\item Dans $H \otimes H'$, l'élément $u * v = b'u + bv - 2u \otimes v$ annule $Q$. On a donc un morphisme de $k$-algèbres $k[t]/(Q) \to H \otimes H'$, $t \mapsto u * v$, comorphisme de $X \times_S X' \to \widetilde S$. En particulier $\operatorname{Tr}(u * v) = bb'$ et $\det(u * v) = b^2c' + b'^2c - 4cc'$.
\item L'involution canonique $\sigma$ de $H$ ($u \mapsto b - u$), étendue par $\sigma \otimes 1$, envoie $u * v$ sur $bb' - u * v$ ; il en va de même de $1 \otimes \sigma'$, d'où $u' * v = u * v'$.
\item Pour deux caractères $\varphi : H \to K$, $\varphi' : H' \to K$ à valeurs dans une $k$-algèbre commutative $K$, posons $z = \varphi(u)$ et $z' = \varphi'(v)$. Alors $(\varphi \otimes \varphi')(u * v) = b'z + bz' - 2zz' = t_{\varphi,\varphi'}(u, v)$, et c'est une racine de $Q$. Pour des traces vérifiant $\operatorname{Tr}1 = 2$, la forme $t_{\varphi,\varphi'}$ est un twist : $t_{\varphi,\varphi'}(x, 1) = \operatorname{Tr}x$ et $t_{\varphi,\varphi'}(1, y) = \operatorname{Tr}y$. Remplacer $\varphi$ par le caractère conjugué $\varphi \circ \sigma$ change $t$ en $bb' - t$.
\end{enumerate}
\end{enonce}

\begin{proof}
(1) On restreint $\det$ à $x_3 = 0$. Pour une forme ternaire $\sum a_ix_i^2 + \sum_{i<j} a_{ij}x_ix_j$, le demi-discriminant est $4a_1a_2a_3 + a_{12}a_{13}a_{23} - a_1a_{23}^2 - a_2a_{13}^2 - a_3a_{12}^2$, et son double est le déterminant de la matrice polaire (règle de Sarrus). Avec $a_1 = 1$, $a_2 = c$, $a_3 = c'$, $a_{12} = b$, $a_{13} = b'$, $a_{23} = bb' - t$, on obtient $4cc' + bb'(bb' - t) - (bb' - t)^2 - cb'^2 - c'b^2 = -Q(t)$. Le discriminant de $Q$ est $(bb')^2 - 4(b^2c' + b'^2c - 4cc') = (b^2 - 4c)(b'^2 - 4c')$, et $Q(bb' - t) = Q(t)$ se développe directement. La matrice de la forme trace est celle de la page~92, avec $t^2 - 2cc'$ en bas à droite ; son déterminant, développé selon la première ligne, vaut $-Q^2$. Comme $Q$ est unitaire de degré~$2$, $k[t]/(Q)$ a pour base $(1, t)$. (2) Dans l'anneau commutatif $H \otimes H'$, avec $U = u \otimes 1$ et $V = 1 \otimes v$, on a l'identité
\[ Q(b'U + bV - 2UV) = (2V - b')^2(U^2 - bU + c) + (b^2 - 4c)(V^2 - b'V + c'), \]
dont le second membre est nul. (3) $(\sigma \otimes 1)(u * v) = b'(b - U) + bV - 2(b - U)V = bb' - (b'U + bV - 2UV)$, et symétriquement. (4) $\varphi \otimes \varphi'$ envoie $U$ sur $z$ et $V$ sur $z'$, d'où la valeur ; l'identité de (2), appliquée dans $K$ aux racines $z$ et $z'$, montre que $b'z + bz' - 2zz'$ annule $Q$. Pour le twist, $t_{\varphi,\varphi'}(x, 1) = 2\varphi(x) + \operatorname{Tr}(x) - 2\varphi(x)$. Enfin $b'(b - z) + bz' - 2(b - z)z' = bb' - (b'z + bz' - 2zz')$.
\end{proof}

\begin{enonce}[Pages 97 et 98, le critère d'Azumaya\piste{83-twist-line-degenerate-locus-is-product}]
Soit $k$ un anneau commutatif quelconque, caractéristique~$2$ comprise, et $(b, c, b', c', t) \in k^5$, sans hypothèse sur $\delta$, $\delta'$. L'algèbre $M = M(b, c; b', c'; t)$ est d'Azumaya sur $k$ (finie, projective, fidèle, et $M \otimes M^{\mathrm{op}} \to \operatorname{End}_k M$, $a \otimes a' \mapsto (x \mapsto axa')$, bijective) si et seulement si $Q(t) = t^2 - bb't + b^2c' + b'^2c - 4cc'$ est inversible dans $k$. De façon équivalente, la forme $\det$ est non dégénérée sur $B = k1 + ku + kv$, c'est-à-dire que $\delta_{P,B}$ est inversible. Sur un corps, $M$ est d'Azumaya si et seulement si $Q(t) \neq 0$ : le lieu non Azumaya de la droite des twists est $\widetilde S = \{Q = 0\}$. Pour $b = c = b' = c' = 0$, $M(0, 0; 0, 0; t)$ est d'Azumaya si et seulement si $t \in k^{*}$ (page~97).
\end{enonce}

\begin{proof}
$M$ est libre de base $e = (1, u, v, uv)$, donc finie, projective et fidèle ; tout se ramène à la bijectivité de $\Phi : M \otimes M^{\mathrm{op}} \to \operatorname{End}_k M$. Dans les bases $(e_i \otimes e_j)$ et $(E_{rs})$, $\Phi$ a pour matrice la matrice $16 \times 16$ dont le coefficient d'indice $((r, s), (i, j))$ est la coordonnée $r$ de $e_ie_se_j$ ; $\Phi$ est bijective si et seulement si son déterminant $D$ est inversible. Un calcul en coordonnées donne $\operatorname{Tr}(x \mapsto axa') = \operatorname{Tr}(a)\operatorname{Tr}(a')$. Composons $\Phi$ avec l'accouplement trace de $\operatorname{End}_k M$, $(f, g) \mapsto \operatorname{Tr}(fg)$, dont la matrice dans la base $(E_{rs})$ est la permutation $(r, s) \mapsto (s, r)$. Comme $\Phi$ est multiplicative, on obtient la forme $(e_i \otimes e_j, e_k \otimes e_l) \mapsto \operatorname{Tr}(e_ke_i)\operatorname{Tr}(e_je_l)$, produit de Kronecker $F \otimes F$ de la forme trace $F = (\operatorname{Tr}(e_ie_j))$. Son déterminant vaut $\det(F)^8 = (-Q^2)^8 = Q^{16}$, d'où $\varepsilon D^2 = Q^{16}$ avec $\varepsilon = \pm 1$ la signature de la permutation. Ainsi $D$ est inversible si et seulement si $Q$ l'est. L'équivalence avec $\delta_{P,B}$ vient de $\delta_{P,B} = -Q(t)$, et le cas nilpotent de $Q = t^2$. L'argument n'utilise que le déterminant de la forme trace, et non sa non-dégénérescence : il vaut donc aussi en caractéristique~$2$, où la remarque de la lecture ne s'appliquait pas.
\end{proof}

\begin{remarque}
Lean (\href{\leanblobmain/Folder83.lean}{\texttt{Folder83.lean}}) démontre le premier énoncé (\texttt{mulM\_assoc}, \texttt{one\_mulM}, \texttt{mulM\_one}, \texttt{mulM\_add}, \texttt{add\_mulM}, \texttt{u\_mul\_u}, \texttt{v\_mul\_v}, \texttt{u\_mul\_v}, \texttt{v\_mul\_u}, \texttt{hamilton\_cayley}, \texttt{nrdM\_mulM}, \texttt{trM\_w}, \texttt{realise\_mulM}, \texttt{realise\_unique}, \texttt{mulM\_inj}), le second (\texttt{nrdM\_B}, \texttt{two\_mul\_halfDisc3}, \texttt{deltaPB\_eq}, \texttt{disc\_Q}, \texttt{Q\_invol}, \texttt{det\_traceForm}, \texttt{Stilde\_rank}, \texttt{Q\_star\_identity}, \texttt{star\_root}, \texttt{starHom}, \texttt{conjL\_star}, \texttt{conjR\_star}, \texttt{productMap\_star}, \texttt{twist\_eq\_productMap}, \texttt{twist\_root}, \texttt{productMap\_star\_root}, \texttt{twist\_one\_right}, \texttt{twist\_one\_left}, \texttt{twist\_conj}) et le troisième (\texttt{isAzumaya\_iff}). Tout vaut sur un anneau commutatif quelconque, caractéristique~$2$ comprise. Le cœur algébrique de la piste est donc démontré en coordonnées. Dans le cas libre, les structures de Hamilton--Cayley forment la droite des $t$, et leur lieu dégénéré $\delta_{P,B} = 0$ est $Q(t) = 0$. Son algèbre $k[t]/(Q)$ reçoit $t \mapsto u * v$, et $Q$ est le polynôme du produit $H * H'$ (cote~84, section~\ref{sec:84}). La structure de torseur n'est démontrée que pour $H$, $H'$ libres, $L$ et $L'$ trivialisés : les torseurs intrinsèques sur une base et la chaîne d'isomorphismes $AZ_P \to \underline{Q}_P \to \underline{Q}_B \to \underline{\Psi}_B \to \underline{\Phi}_B \to \underline{L}_P$ ne sont pas formalisés. Le critère \q{Azumaya si et seulement si $Q$ inversible} de la page~98, qui fait de $Q = 0$ le lieu non Azumaya, est démontré pour la notion d'algèbre d'Azumaya de mathlib (\texttt{IsAzumaya}), sur l'algèbre \texttt{MAlg} qui munit $k^4$ de la table de la page~92 : \texttt{isAzumaya\_iff}, par \texttt{mulLeftRight\_coords}, \texttt{sum\_diag\_mul}, \texttt{sandwich\_gram}, \texttt{sandwich\_det\_sq} et \texttt{isUnit\_sandwichMatrix\_iff}, avec les variantes \texttt{isAzumaya\_iff\_deltaPB}, \texttt{isAzumaya\_iff\_ne\_zero} et \texttt{isAzumaya\_nilpotent\_iff} (page~97). Sa forme intrinsèque, pour deux sous-algèbres de rang~$2$ transverses d'une algèbre localement libre de rang~$4$, ne l'est pas. Pour l'identification $\widetilde S \simeq X \wedge X'$, quotient de $X \times X'$ par $\sigma \times \sigma'$ quand $(2, \delta, \delta')$ engendrent l'idéal unité, seul est démontré le cœur algébrique, les invariants de $\sigma \otimes \sigma'$ égaux à $k[u * v]$ si $(2, b, b')$ est l'idéal unité, et il l'est dans \texttt{Folder84.lean} (\texttt{invariants\_of\_span\_eq\_top}). Les traces de $H$ et $H'$ sont des hypothèses ($\operatorname{Tr}1 = 2$, $\operatorname{Tr}u = b$), et non la trace de mathlib.
\end{remarque}
